\bookchapter{第二十三章附：GitHub入門指南——從代碼託管到協作平臺}

\chapterintro{上章吾輩講了“編程之思”，而真要將編程之思落地、與他人共做項目，即繞不開一名——GitHub。無論君為計算機系學生要參賽、欲存己之代碼作品，抑或單純欲觀全球開發者折騰何物，GitHub皆君必須認識之平臺。此節吾輩自最基礎之概念講起，助君建一清晰之認知。}

\booksection{一、GitHub究竟爲何物？}

簡單言之，GitHub乃一基於Git之\textbf{代碼託管與協作平臺}。其於2008年上線，2018年爲微軟所收購，至今仍為全球最大之開源代碼託管平臺，託管着數億個倉庫，覆蓋幾乎所有主流編程語言。

若以一類比理解之：\textbf{Git乃君電腦上之日記本，GitHub則一公開之博客平臺}——君可發佈己之作品，關注他人之項目，參與討論與協作。其不僅為“程序員之後臺倉庫”，更似全球開發者共建之一座“技術城”——項目於此誕生、迭代、被討論、被複用。

於普通人而言，GitHub亦非遙不可及。多數非程序員會用其託管筆記、電子書、個人博客，甚至用以備份配置文件。只要君有一絲“欲將己之事物整理好並分享出去”之需，GitHub皆值得君瞭解。

\booksection{二、Git與GitHub之別（關鍵概念！）}

多人初觸時常混此兩詞，此處爲君梳理清楚：

\begin{center}
\begin{tabularx}{\linewidth}{|X|X|X|}
\hline
\textbf{對比項} & \textbf{Git} & \textbf{GitHub} \tabularnewline
\hline
本質 & 版本控制工具（軟件） & 代碼託管平臺（網站） \tabularnewline
運行位置 & 本地電腦上 & 雲端（瀏覽器訪問） \tabularnewline
核心職責 & 記錄文件之每次變更（何以改、誰改、何時改） & 將Git倉庫放至雲端，方便存儲、展示、協作 \tabularnewline
為否需網絡 & 不需，離線可用 & 需聯網訪問 \tabularnewline
歸屬 & 開源軟件，由Linus Torvalds（Linux之父）創建 & 微軟旗下產品 \tabularnewline
同類競品 & SVN、Mercurial & GitLab、Gitee（碼雲）、Bitbucket \tabularnewline
\hline
\end{tabularx}
\end{center}

簡言之：\textbf{Git管“何以記錄變更”，GitHub管“變更放於何處、何以協作”}。Git乃工具，於本地運行；GitHub乃平臺，將Git倉庫放至雲端，方便協作與展示。君可只用Git不用GitHub（譬如自己寫代碼做版本管理），然用GitHub必用Git（因GitHub即基於Git）。

\booksection{三、GitHub之核心概念與功能}

\booksubsection{基礎概念}

動手之前，先識數個高頻詞彙，否則觀GitHub之界面將一頭霧水：

\begin{itemize}
  \item \textbf{倉庫（Repository，簡稱Repo）}：一個項目對應一個倉庫，內中存放此項目之所有代碼文件、歷史版本與文檔。可理解爲一"項目文件夾"。
  \item \textbf{提交（Commit）}：對代碼之一次修改記錄。每次Commit皆附說明信息，記"此次改了何、何故改"。此乃Git/GitHub能"時光倒流"之基礎。
  \item \textbf{分支（Branch）}：自主線上分出之獨立開發線。譬如君欲加一新功能，可開一分支單獨開發，開發畢再合併回主線，不影響他人。
  \item \textbf{克隆（Clone）}：將雲端倉庫完整複製一份至本地電腦，此後即可於本地修改代碼。
  \item \textbf{推送（Push）/拉取（Pull）}：Push乃將本地代碼上傳至GitHub，Pull乃自GitHub下載最新代碼至本地。此乃團隊協作時同步代碼之兩個核心動作。
  \item \textbf{合併請求（Pull Request，簡稱PR）}：當君於一分支上改完代碼，欲合併至主分支時，發起一PR，令團隊成員審查代碼後再合併。此乃GitHub協作之核心機制。
  \item \textbf{議題（Issue）}：用以記錄bug、功能需求、討論話題。每個Issue如一討論帖，可指派負責人、打標籤。
  \item \textbf{Fork}：將他人之倉庫複製一份至己之賬號下，君可於己之副本上隨意修改，不影響原倉庫。改好後可予原作者發PR，請求合併君之改進。
\end{itemize}

\booksubsection{核心功能}

GitHub所供不止代碼存儲，尚包括：

\begin{itemize}
  \item \textbf{版本控制}：記錄每次代碼更改之歷史，改壞代碼隨時回至上一正常版本。
  \item \textbf{團隊協作}：支持多用戶對同一項目之並行開發，多人同時開發不會覆蓋、不會衝突。
  \item \textbf{問題跟蹤}：用Issue記錄bug與功能請求，方便團隊跟進。
  \item \textbf{項目管理}：集成了待辦事項、項目看板（Projects）等管理工具，相當於內置一輕量級Trello。
  \item \textbf{自動化}：經\textbf{GitHub Actions}實現代碼提交後自動編譯、測試、打包、部署，乃現代"CI/CD（持續集成/持續部署）"之核心工具。
  \item \textbf{文檔管理}：README介紹項目，Wiki存放詳細架構與開發規範，Pages甚至可徑託管靜態網站。
  \item \textbf{代碼託管與展示}：每個倉庫皆有漂亮之網頁展示，支持語法高亮，乃君技術作品之"門面"。
  \item \textbf{開源社區}：君可Star（藏）喜歡之項目、Watch（關注）項目動態、Follow（關注）開發者，參與全球開源協作。
\end{itemize}

\booksection{四、註冊與上手：五分鐘創建君之第一個倉庫}

\booksubsection{註冊賬號}

\begin{enumerate}
  \item 訪問\href{https://github.com}{github.com}，點擊"Sign up"。
  \item 填寫郵箱、設置密碼與用戶名（用戶名將出現於君之個人主頁地址中，建議用英文，譬如\texttt{github.\allowbreak com/\allowbreak yourname}）。
  \item 完成郵箱驗證與人機驗證。
  \item 選免費方案（Free）即可，對個人用戶完全夠用。
\end{enumerate}

\booksubsection{創建第一個倉庫}

\begin{enumerate}
  \item 登錄後點擊右上角"+"號，選"New repository"。
  \item 填寫倉庫名（Repository name），譬如\texttt{my-\allowbreak first-\allowbreak project}。
  \item 寫一段簡短描述（Description），令他人知此倉庫爲何用。
  \item 選公開（Public）或私有（Private）。公開倉庫所有人可見，私有唯己與協作者可見。免費賬號亦支持無限私有倉庫。
  \item 勾選"Add a README file"——README乃項目之"說明書"，每個倉庫皆建議有之。
  \item 點擊"Create repository"，君之第一個GitHub倉庫即誕生矣！
\end{enumerate}

\begin{tcbinfo}{小知識：README爲何物}
README.md乃一Markdown格式之文件，乃他人打開君倉庫時第一眼所見之內容。一個好之README通常含：項目簡介、安裝方法、使用示例、許可證、聯繫方式。寫好README乃程序員之基本素養，亦項目能否爲他人所用之第一道門檻。

\end{tcbinfo}

\booksection{五、基本工作流：自本地至雲端}

理解GitHub之工作流，核心乃記此循環：\textbf{本地修改 暫存 提交 推送 雲端更新}。下以大白話釋每一步：

\begin{enumerate}
  \item \textbf{修改文件}：於本地電腦上以編輯器寫代碼或改文檔。
  \item \textbf{git add}：將修改之文件加入"暫存區"，相當於告Git"此等改動吾準備記錄下來"。
  \item \textbf{git commit}：將暫存區之改動正式記成一"提交"，附上一句說明，譬如"修復登錄bug"。此步唯於本地完成。
  \item \textbf{git push}：將本地之提交上傳至GitHub雲端，如此他人即能見之。
  \item \textbf{git pull}：於始工作前，先自雲端拉取最新代碼，確保本地與團隊其他人保持同步。
\end{enumerate}

於新手而言，記此口訣：\textbf{“先pull後改，改完add，commit帶說明，最後push上去”}。養此習慣，團隊協作即不會亂。

\booksection{六、團隊協作流程：Pull Request何以運作}

GitHub最強之處於其協作機制。一個典型之團隊協作流程乃：

\begin{enumerate}
  \item \textbf{開分支}：自主分支（通常叫main或master）開一新分支，譬如\texttt{feature/\allowbreak login}，於此分支上開發新功能。
  \item \textbf{開發並提交}：於新分支上寫代碼，多次commit記錄進度。
  \item \textbf{推送至GitHub}：將分支push至雲端倉庫。
  \item \textbf{發起Pull Request}：於GitHub網頁上發起PR，請求將君之分支合併至主分支。寫清楚此PR做了何改動、何故改。
  \item \textbf{代碼審查（Code Review）}：團隊成員查看君之代碼，提出意見或建議修改。君可根據反饋繼續commit，PR將自動更新。
  \item \textbf{合併}：審查通過後，點擊"Merge"將代碼合併至主分支。君之貢獻即正式進入項目矣！
\end{enumerate}

此套流程保證了\textbf{任何代碼進入主分支前皆經審查}，大大降低出錯概率。開源項目亦用此套機制接受全世界陌生人之貢獻——君完全可於某知名開源項目上提交PR，被採納後君即此項目之貢獻者矣。

\booksection{七、GitHub之典型應用場景}

\booksubsection{學生參賽與課程項目}

計算機系學生參加ACM、黑客松、畢業設計，幾乎皆離不開GitHub。其既能作團隊協作平臺，又能作作品展示窗口——比賽評委可徑觀君之倉庫，瞭解代碼質量與協作過程。多數公司招聘時亦會要求附上GitHub主頁，作爲技術能力之證明。

\booksubsection{個人知識庫與筆記託管}

非程序員亦能用GitHub。多人用其託管學習筆記、技術博客（配GitHub Pages可免費建站）、電子書。Markdown格式寫作 + Git版本管理 + 雲端備份，三合一。

\booksubsection{參與開源項目}

GitHub乃全球開源運動之核心陣地。君可Fork喜歡之項目，修改後發PR回饋社區；亦可單純Star藏好項目，跟蹤更新。參與開源乃提升技術、積累聲譽之絕佳途徑。

\booksubsection{自動化部署與CI/CD}

借GitHub Actions，可實現“代碼一push，自動測試+自動部署至服務器”之全自動流程。多數個人博客、文檔站點皆如此搭建，寫完文章push之，網站即自動更新矣。

\booksubsection{代碼備份與多設備同步}

將代碼倉庫當作雲盤用，換電腦、多設備辦公時，git pull之即能拿到最新代碼，再不用U盤拷來拷去。

\booksection{八、予新手之數點建議}

\begin{itemize}
  \item \textbf{提交說明要寫清楚}：勿寫"update"此等廢話說明。好之commit message寫"修復用戶登錄後跳轉錯誤頁面之問題"，半年後回觀亦知當時幹了啥。
  \item \textbf{勤提交，小步快跑}：每完成一個小功能即commit一次，而非寫了一整日方提交一次。如此出問題更易定位。
  \item \textbf{用分支開發，勿徑改主分支}：哪怕為個人項目，養成開分支之習慣，未來團隊協作時即不會手忙腳亂。
  \item \textbf{善用.gitignore}：此文件用以告Git哪些文件勿納入版本管理（譬如密碼配置、編譯產物、依賴文件夾），避免將垃圾文件傳上去。
  \item \textbf{多觀優秀開源項目}：逛逛GitHub Trending，觀高星項目之代碼結構、README寫法、協作流程，乃提升最快之免費方式。
  \item \textbf{申請GitHub學生包}：學生認證後可免費獲得大量付費工具與服務之額度（Copilot、域名、雲服務等），價值上千美元。
\end{itemize}

\begin{tcbwarn}{安全提醒}
\begin{itemize}
  \item 千萬勿將密碼、API密鑰、數據庫連接串等敏感信息提交至GitHub，尤以公開倉庫爲甚——一旦泄露，將爲爬蟲秒抓，後果嚴重。
  \item 不小心提交了敏感信息，光刪除文件不夠（歷史記錄裏尚有），需用\texttt{git filter-\allowbreak branch}或專業工具清理歷史，並立即更換泄露之密鑰。
  \item 私有倉庫亦非絕對安全，敏感信息最好用環境變量管理，不進代碼倉庫。
\end{itemize}

\end{tcbwarn}

\begin{tcbinfo}{國內訪問GitHub之注意事項}
GitHub於國內訪問有時不穩，或出現打不開、克隆慢之情況。可嘗試：(1) 用國內鏡像站（如GitClone、ghproxy）加速克隆；(2) 配置Hosts或用代理；(3) 若主要面向國內協作，亦可考慮國產替代品\textbf{Gitee（碼雲）}，功能類似且訪問速快。多數團隊會同時於GitHub與Gitee上託管，兼顧國際化與國內速度。

\end{tcbinfo}

\textbf{進階玩法：以GitHub Pages搭免費個人博客}\par

GitHub Pages允許君免費託管一靜態網站，配Jekyll、Hexo、Hugo等靜態站點生成器，可輕鬆搭建個人博客。寫畢Markdown文章push至倉庫，博客自動更新，且自帶HTTPS與自定義域名。無數技術博主之博客皆如此搭成，零成本、無廣告、版本可控。

\booksection{九、總結：何故每人皆該瞭解GitHub}

GitHub早已不只是“程序員之工具”，其更像是一\textbf{數字時代之工作方式}數字時代之工作方式

對計算機系之學生而言，GitHub更是必備技能：它是君之作品集、協作平臺、學習資源庫，亦乃未來求職時最有說服力之“簡歷”。早點註冊賬號，將君之課程作業、小項目傳上去，慢慢積累，數年後君將擁有一份令面試官眼前一亮之技術履歷。

\begin{tcbpractice}{試爲之：完成君之GitHub初體驗}
\begin{enumerate}
  \item 註冊一GitHub賬號，完善個人資料（頭像、簡介、感興趣之技術領域）。
  \item 創建一名爲\texttt{hello-\allowbreak github}之公開倉庫，勾選"Add a README file"。
  \item 進入倉庫，點擊README.md之編輯按鈕（鉛筆圖標），用Markdown寫一段自我介紹，提交更改。
  \item 於GitHub搜索欄搜一君感興趣之關鍵詞（譬如\texttt{python-\allowbreak tutorial}），Star兩個君喜歡之項目。
  \item （進階）於電腦上安裝Git，將倉庫Clone至本地，修改README後push回雲端，完成一次完整之"本地雲端"工作流。
\end{enumerate}

\end{tcbpractice}
